\documentclass[11pt]{article}
\usepackage[margin=1in]{geometry}
\usepackage[T1]{fontenc}
\usepackage[utf8]{inputenc}
\usepackage{lmodern}
\usepackage{microtype}
\usepackage{amsmath,amssymb,amsthm,mathtools}
\usepackage{booktabs}
\usepackage{array}
\usepackage{longtable}
\usepackage{graphicx}
\usepackage{xcolor}
\usepackage{hyperref}
\usepackage[nameinlink,noabbrev]{cleveref}
\usepackage{enumitem}
\usepackage{listings}
\usepackage{csquotes}

\hypersetup{
    colorlinks=true,
    linkcolor=blue,
    citecolor=blue,
    urlcolor=blue
}

\setlist[itemize]{leftmargin=1.5em}
\setlist[enumerate]{leftmargin=1.75em}

\newcommand{\trellis}{\textsc{Trellis}}
\newcommand{\platformrequest}{\texttt{PlatformRequest}}
\newcommand{\semanticcontract}{\texttt{SemanticContract}}
\newcommand{\valuationcontext}{\texttt{ValuationContext}}
\newcommand{\productir}{\texttt{ProductIR}}
\newcommand{\eventprogramir}{\texttt{EventProgramIR}}
\newcommand{\controlprogramir}{\texttt{ControlProgramIR}}
\newcommand{\familyir}{family IR}
\newcommand{\code}[1]{\texttt{#1}}


\title{Trellis I: A Request-to-Outcome Compiler for Agent-Assisted Derivatives Pricing}
\author{Draft}
\date{\today}

\begin{document}
\maketitle

\begin{abstract}
This paper should explain how \trellis{} turns ambiguous pricing requests into
defensible deterministic pricing outcomes. The central claim is that
agent-pricing becomes viable when the LLM is constrained to semantic parsing,
planning, and guarded artifact generation around a typed compiler and a
deterministic numerical runtime, rather than being placed in the pricing hot
path.
\end{abstract}

\tableofcontents

\section{Introduction}

\subsection{Problem}
Most pricing libraries are organized around instruments and pricing engines.
This is true whether the implementation is explicitly engine-based or more
product-centered. A European option has one pricing interface, a Bermudan
swaption another, a callable bond another, and each one is connected to the
pricing methods that the library already knows how to apply to that product.

This design is sensible. It is reliable, easy to explain, and usually the
right way to ship a production library. But it also means that similar
numerical structure is often hidden behind different product interfaces. A
callable bond and a Bermudan swaption are different products, yet both require
an event schedule, an exercise rule, and a backward-induction style numerical
method. Likewise, a European vanilla option might be priced analytically, by a
lattice, by a PDE, or by Monte Carlo, but in a conventional library these are
still usually introduced through instrument-specific pricing contracts.

This product-centered organization is easy to see from ordinary library usage.
In a FinancePy-style workflow, valuation is exposed at the product object:

\begin{lstlisting}[language=Python]
# Representative FinancePy-style workflow
model = BlackScholes(...)
product = EquityAmericanOption(...)
price = product.value(valuation_date, market_data, model)
\end{lstlisting}

In a QuantLib-style workflow, the instrument and engine are separated, but the
engine is still attached through an instrument-specific contract:

\begin{lstlisting}[language=Python]
# Representative QuantLib-style workflow
payoff = ql.PlainVanillaPayoff(ql.Option.Call, strike)
exercise = ql.EuropeanExercise(maturity)
option = ql.VanillaOption(payoff, exercise)

process = ql.BlackScholesMertonProcess(...)
engine = ql.AnalyticEuropeanEngine(process)
option.setPricingEngine(engine)
price = option.NPV()
\end{lstlisting}

Both workflows are clear and practical. They also make the usual abstraction
boundary visible. To move from a European vanilla option to a barrier option,
an Asian option, a Bermudan swaption, or a callable bond, the user does not
usually keep one generic claim representation and only change control, state,
and timeline metadata. Instead, the user typically moves to a different
product class, and often to a different family of compatible pricing engines.

That is the design point where \trellis{} begins. The central question is not
simply ``which pricer belongs to this product class?'' The earlier question is:
what is the product, in computational terms? Is it a terminal payoff or a
schedule-driven claim? Does it have holder exercise, issuer exercise, or no
exercise at all? Is it terminal-state, path-dependent, or event-dependent? What
market objects and model bindings are required before a numerical method can be
applied?

In the shipped system, those questions are made explicit before method
selection. The compiler does not begin from a raw ``(product, method)'' pair.
It begins from typed product semantics and separates:
\begin{itemize}
    \item contract meaning
    \item valuation context and market binding
    \item admissible numerical families
    \item exact checked backend bindings
\end{itemize}

Concretely, the semantic layer records objects such as payoff family,
exercise style, schedule dependence, state dependence, controller semantics,
and event structure. Those are then combined with a separate valuation layer
carrying model, measure, discounting, market-data, and reporting policy. Only
after that does the compiler lower the request onto a bounded numerical family
such as \code{AnalyticalBlack76IR}, \code{TransformPricingIR},
\code{ExerciseLatticeIR}, or \code{EventAwareMonteCarloIR}.

The point is not a universal solver and it is not free-form factorization for
its own sake. The point is narrower: many products can share the same
computational lane once their control, state, and timeline obligations are made
explicit. In that sense, Trellis tries to make semantic structure, rather than
product-local route identity, the primary source of truth.

This distinction matters for agent-assisted pricing. A prompt or structured
request should not cause the system to improvise a new product-local pricing
class whenever it encounters a new surface description. The better question is
whether the request can be compiled into an already-supported semantic family
and then lowered onto an already-governed deterministic lane.

\subsection{Contribution}
The contribution of this paper is a request-to-outcome architecture for
agent-assisted pricing built around four ideas:
\begin{itemize}
    \item unified request normalization
    \item typed semantic compilation
    \item admissibility and lowering before pricing
    \item deterministic validation and traceability
\end{itemize}

The more specific claim is a change in abstraction boundary. Instead of asking
which product class owns a pricing method, Trellis asks which typed semantic
family can be admitted to which computational lane under which market and
control obligations. In the current shipped architecture, route identifiers are
increasingly reduced to backend-binding and compatibility roles, while semantic
contracts, family registries, and bounded family IRs carry the constructive
meaning.

\section{From Front Doors to a Canonical Request}

Cover:
\begin{itemize}
    \item \code{ask(...)}
    \item \code{Session}
    \item \code{Pipeline}
    \item structured user-defined products
    \item comparison and task-runner flows
\end{itemize}

Key objects:
\begin{itemize}
    \item \platformrequest{}
    \item \code{CompiledPlatformRequest}
    \item \code{ExecutionPlan}
\end{itemize}

\section{Semantic Compilation}

Cover:
\begin{itemize}
    \item \semanticcontract{}
    \item semantic validation
    \item \valuationcontext{}
    \item required data and market binding
    \item \productir{}
\end{itemize}

\section{Admissibility, Lowering, and Route Selection}

Cover:
\begin{itemize}
    \item route registry and build gate
    \item \eventprogramir{} and \controlprogramir{}
    \item lowering onto bounded family IRs
    \item helper-backed deterministic routes
\end{itemize}

\section{Validation and Traces}

Cover:
\begin{itemize}
    \item validation as a compiled contract
    \item deterministic checks versus residual model review
    \item canonical trace surfaces and replayability
\end{itemize}

\section{Case Studies}

Suggested cases:
\begin{itemize}
    \item one direct deterministic pricing request
    \item one guarded build path
    \item one comparison-task path
\end{itemize}

\section{What This Paper Does Not Claim}

\begin{itemize}
    \item the LLM does not perform deterministic pricing
    \item the compiler does not imply universal numerical support
    \item unsupported capability remains bounded by documented limitations
\end{itemize}

\section{Conclusion}

Restate the main result: the meaningful innovation is a governed semantic
compiler around deterministic quantitative engines.

\end{document}
